% Generated draft: requires mathematical and source review.
\documentclass[11pt]{article}
\usepackage[margin=25mm]{geometry}
\usepackage{fontspec}
\setmainfont{Latin Modern Roman}
\usepackage{amsmath,amssymb,mathtools,unicode-math}
\setmathfont{Latin Modern Math}
\usepackage{enumitem,needspace,xurl,hyperref}
\setlength{\parindent}{0pt}
\setlength{\parskip}{5pt}
\setlength{\emergencystretch}{3em}
\newcommand{\N}{\mathbb{N}}
\newcommand{\Z}{\mathbb{Z}}
\newcommand{\Q}{\mathbb{Q}}
\newcommand{\R}{\mathbb{R}}
\newcommand{\C}{\mathbb{C}}
\newcommand{\F}{\mathbb{F}}
\newcommand{\A}{\mathbb{A}}
\newcommand{\PP}{\mathbb{P}}
\newcommand{\E}{\mathbb{E}}
\newcommand{\eps}{\varepsilon}
\newcommand{\dd}{\,\mathrm{d}}
\begin{document}
\section*{UnsolvedMath ID: 2546}
\textbf{Source problem code:} catalogue entry
\subsection{Definitions and mathematical statement}
A tdlc group is a Hausdorff topological group that is totally disconnected and locally compact. A group \(G\) is soluble if its derived series \(G^{(0)}=G\), \(G^{(j+1)}=[G^{(j)},G^{(j)}]\) reaches the trivial group after finitely many steps.

The supplied definition calls a tdlc group constructible if it is obtained starting from the trivial group through a sequence of profinite extensions and ascending HNN-extensions. A profinite group is an inverse limit of finite discrete groups. For reference, an algebraic ascending HNN-extension of a group \(H\) along an injective endomorphism \(\varphi:H\to H\) has presentation
\[
\langle H,t\mid t^{-1}ht=\varphi(h)\ (h\in H)\rangle.
\]
The topological version additionally needs a specified tdlc topology and continuity/openness hypotheses, which the supplied short definition does not state. Likewise, an extension of \(H\) by a profinite kernel \(K\) has an exact sequence \(1\to K\to E\to H\to1\), but the source does not explicitly say whether that kernel convention or a profinite quotient convention is intended by profinite extension, nor whether the construction sequence must be finite.

In a specified module category, type \(FP_\infty\) means that the trivial coefficient module has a projective resolution
\[
\cdots\longrightarrow P_2\longrightarrow P_1\longrightarrow P_0
\longrightarrow R\longrightarrow0
\]
with every \(P_j\) finitely generated. For tdlc groups one must declare the coefficient ring \(R\) and module category. One usual setting uses rational discrete modules: \(R=\Q\), each module is discrete and its \(G\)-action continuous, equivalently each vector has open stabilizer. The supplied record does not explicitly specify that convention, so it is presented as a possible framework, not silently fixed as the only meaning of \(FP_\infty\).

Subject to the intended tdlc construction and homological conventions, the question is the converse implication
\[
\forall G\text{ a soluble tdlc group},\qquad
G\text{ of type }FP_\infty\Longrightarrow G\text{ constructible}.
\]
It does not ask merely for the forward implication. The background, attributed to G. C. Cook, reports that soluble constructible tdlc groups have type \(FP_\infty\), citing Cook--Castellano (2020). The supplied literature triage dated 2026-08-17 labels the row partially solved because of this known direction and leaves the converse characterization open. Both the registry and the triage retain open status for the displayed converse. No unmentioned coefficient choice, HNN topology or extension direction is imposed to manufacture a fully specified theorem.
\subsection{Short English statement}
Must a soluble totally disconnected locally compact group with the homological finiteness property \(FP_\infty\) be built from the trivial group by profinite and ascending HNN extensions? The supplied forward direction is known; the converse is open in its assessment, and the exact homological and topological construction conventions need to be taken from the cited source.
\subsection{Sources}
\begin{itemize}
\item Local UnsolvedMath registry, permanent ID 2546, source code catalogue entry, ``Kourovka Notebook Problem 21.37,'' from \emph{Kourovka Notebook}, New Problems, 21st issue (2026), original problem 21.37; attributed to G. C. Cook. Includes supplied literature triage dated 2026-08-17.
\item G. C. Cook and I. Castellano, \emph{Journal of Algebra} 543 (2020), 54--97; supplied reference for the forward constructible-to-\(FP_\infty\) result. No paper title, DOI or URL is supplied.
\item \emph{Kourovka Notebook}, 21st issue (2026), Problem 21.37, supplied maintained source: \url{https://www.kourovka-notebook.org/}.
\end{itemize}
\end{document}
